%% IEEE conference paper — compile with: pdflatex planning_failover.tex (x2)
%% Requires the IEEEtran document class (ships with TeX Live; Overleaf-ready).
%% Author block, affiliation, and reference details are placeholders — fill before submission.
\documentclass[conference]{IEEEtran}

\usepackage{cite}
\usepackage{amsmath,amssymb}
\usepackage{graphicx}
\usepackage{xcolor}
\usepackage{url}
\usepackage[T1]{fontenc}

\begin{document}

\title{Evidence-Driven, Failure-Resilient Planning for\\
Multi-Agent Software-Engineering Automation}

\author{
\IEEEauthorblockN{Author One, Author Two}
\IEEEauthorblockA{\textit{Affiliation / Organization} \\
City, Country \\
\texttt{\{author\}@example.com}}
}

\maketitle

\begin{abstract}
Many systems now use a large language model (LLM) to \emph{plan} software-
engineering work: the user writes a request in plain language, and the model
chooses which specialized agents should run, and in what order, to analyze,
generate, test, patch, and report on code. In practice this planning step is the
part that breaks most often. The model invents agent names that do not exist,
builds a plan for a request that is unclear or not even about software, drops
needed steps from large plans without warning, and stops the whole workflow when
the model is down or rate limited. We present a planning subsystem that handles
these problems with three parts that work together. First, an
\emph{evidence-driven plan stage} collects cheap, rule-based facts about the
request---what the user is asking for, what the repository actually contains, and
which online agents fit---before it makes a single model call, and then has a
second model call check the plan. Second, a set of \emph{rule-based ``reflection''
checks} fix the plan before it runs, drop steps that later turn out to be
pointless, and make the final report honest (for example, it never reports a
passing build when tests failed). Third, a \emph{fallback layer} keeps the system
running when something fails: model calls fall back from a hosted model to a local
one to a safe default, every fact-gathering step is allowed to fail without
blocking the request, and the system pauses to ask a human when the request is
unclear or the plan looks risky. We describe the design, give a simple way to test
planning on large and messy requests, and report real problems this testing
found---a hidden bug that silently dropped the final report, and a model rate limit
that capped how fast we could plan.
\end{abstract}

\begin{IEEEkeywords}
LLM agents, multi-agent systems, workflow orchestration, human-in-the-loop, fault
tolerance, software engineering automation.
\end{IEEEkeywords}

\section{Introduction}
A common way to automate software work is to break a user request into a series of
specialized agents or tools. A \emph{planner}---usually an LLM---reads the request
and decides the steps: analyze the repository, find bugs, write tests, patch,
re-test, run a security check, and write a report. The quality of the whole run
depends on this plan.

But the planner is where most failures start. We see four problems again and
again: (1)~\textbf{broken output}---small or local models return invalid JSON,
repeated or non-existent agent names, or far too many steps;
(2)~\textbf{wrong choice when the request is unclear}---a vague, contradictory, or
off-topic request gets a made-up plan instead of a question back to the user;
(3)~\textbf{silently losing steps}---length limits, first added to protect against
model mistakes, quietly delete needed steps from large plans, most often the final
report; (4)~\textbf{one point of failure}---when the planner model is down, slow,
or rate limited, the workflow stops or falls back to a poor default.

A bigger model does not fix all of this. A stronger model helps with problems (1)
and (2), but does nothing for (3) and (4)---and even a strong model is an outside
service with a limited request rate.

We describe the planning part of an open multi-agent platform (the
``orchestrator'') built on a state-machine workflow engine. Our three
contributions, treated as equally important, are: an \textbf{evidence-driven plan
stage} that first gathers three kinds of cheap, rule-based facts---request intent,
a quick repository scan, and a fit score for each online agent---before any model
call, and then has a second model call (a ``critic'') review the plan and list the
agents it left out; \textbf{rule-based reflection}---three small checks (no model
call) that fix the plan before it runs, remove steps later shown to be pointless,
and make the final status honest; and a \textbf{fallback layer}---a
hosted~$\rightarrow$ local~$\rightarrow$ default model chain, fact steps that fail
safely, human checkpoints triggered by unclear requests or risky plans, and a side
channel for non-blocking warnings. We also give a \textbf{simple testing method}
for large and unclear requests that checks \emph{behavior} (Did
it ask for clarification? Are the required agents there? Is the report kept?)
instead of comparing the plan to one ``correct'' answer.

\section{Related Work}
\textbf{LLM agents and tool use.} Step-by-step reasoning~\cite{cot} improves model
answers; ReAct~\cite{react} mixes reasoning and actions; Toolformer~\cite{toolformer}
teaches a model to call tools; open-ended agents such as Voyager~\cite{voyager}
chain model calls freely. Our planner is on purpose \emph{not} open-ended: it
produces a fixed list of steps that a plain engine then runs, so the result is easy
to read and check.

\textbf{Coordinating many agents.} Frameworks like AutoGen~\cite{autogen} and
graph-based engines~\cite{langgraph} run several agents together. We use such an
engine but focus on the \emph{planning} step and how to keep it working---which
these frameworks leave to the application.

\textbf{Self-checking and ``model-as-judge.''} Reflexion~\cite{reflexion} has a
model critique and revise its own output; ``LLM-as-a-judge''~\cite{judge} uses one
model to score another. Our plan critic is a judge specialized to picking the right
steps, but it only runs on confident plans, and its verdict either auto-fixes the
plan or sends it to a human.

\textbf{Routing by capability.} Service systems route work by what each component
advertises it can do. We use a simple, rule-based score over what each online agent
advertises, rather than a model or a learned router, so routing stays cheap and new
agents become usable as soon as they come online.

\textbf{Asking a human, and clean output.} Asking a human before risky actions is
standard, and constrained decoding helps a model return valid output. We use both:
the human checkpoints are triggered automatically by signals the plan stage
computes, and we keep a forgiving JSON reader as a backup behind the stronger
hosted model.

\section{System Architecture}
The orchestrator runs a set of specialized \textbf{KIO agents}. Each is a small
network service that does one thing: analyze a repository, find bugs, generate
tests, generate patches, re-run tests, generate new code, run a security/OWASP
check, optimize a model for low-power devices, build extra automated tests, or
write the final report. A workflow is just an ordered list of these agents.

The engine is a state machine---a graph of steps with saved state, so a run
survives a restart. The first step is \texttt{plan}; it turns the request (plus any
attached code or repository) into a checked list of agents. From \texttt{plan}, a
small router sends the run to one of three places: \texttt{plan\_clarify} (too
unclear to plan; ask the user and loop back), \texttt{plan\_review} (risky plan;
pause for a human to approve), or \texttt{run\_kio} (run the agents). The rest of
the paper is about the \texttt{plan} step and the machinery that keeps it reliable.

\section{The Evidence-Driven Plan Stage}
\label{sec:plan}
Instead of asking the model to choose steps from the raw request, we first build a
short set of \emph{facts}, then make one hosted model call, then check the result
(Fig.~\ref{fig:stage}).

\begin{figure}[t]
\centering
\footnotesize
\begin{verbatim}
        +------------- plan node -------------+
 req -> | intent -+                           |
(+ctx)  | scan   -+- facts - LLM planner ----+|
        | bids   -+  block   (+left-out list)||
        |                          critic <--+|
        |       clarification gate            |
        +----+--------------+-------------+----+
             v              v             v
       plan_clarify    plan_review    run_kio
        (ask user)   (risky plan,    (run agents)
            |        human approves)
            +-> loop back to plan with the answer
\end{verbatim}
\caption{The plan stage: cheap rule-based facts feed one hosted model call, which a
second model call (the critic) checks, then the run goes to clarification, human
review, or execution.}
\label{fig:stage}
\end{figure}

\textbf{A. Reading the request intent.} A rule-based step (no model call) turns the
request into a few simple labels: the \texttt{task\_type} (security, optimization,
testing, bug fix, code generation, analysis, or unknown), whether code is already
provided, and a \texttt{risk\_level}. It uses word patterns plus facts we already
have, so it is fast and free, and it makes routing easy to follow (a security or
high-risk request should include the security agent).

\textbf{B. A quick repository scan.} When the request points at an existing
repository, we scan file names only (we do not read file contents) to produce a
short fact sheet: which languages are present, whether there are tests, which
dependency files exist, and whether there are security-sensitive areas (login,
crypto, secrets). This bases the choice on what the code \emph{is}, not on the
wording of the request.

\textbf{C. Agent fit scores (``bidding'').} Each agent advertises what it can do
when it comes online. For a request, we score every \emph{online} agent's
advertised text against the request by simple word overlap and pass the top matches
to the planner. Because the score uses text we already have, there is no extra round
trip, and an agent added after deployment becomes usable without changing the
planner prompt.

\textbf{D. The model planner and the ``left-out'' list.} The three fact blocks, plus
a list of what each agent does, go to one hosted model call. It returns the list of
steps, a one-line reason, and a list of agents it \emph{considered but left out},
each with a reason. That left-out list is given to the critic so a wrong omission
can be caught, and it makes the plan easy to audit. The output is checked against
the valid agent names, de-duplicated, length-limited (Section~\ref{sec:resil}-C), and read
by a forgiving JSON reader as a backup.

\textbf{E. The plan critic (only on confident plans).} On confident plans only (a
fallback plan already goes to a human), a second model call reviews the plan against
the request and the same facts. It returns \texttt{ok} (keep it), \texttt{adjust}
(use a corrected list), or \texttt{uncertain} (send it to human review). The facts
improve what the planner \emph{sees}; the critic checks what the planner
\emph{produced}, where the actual mistake is visible.

\textbf{F. The clarification gate.} Before planning, a quick check decides whether
the request is a clear software task. If it is not---nonsense, off-topic, empty, or
too vague---the stage does not guess; it asks the user one question with up to four
concrete options plus free text, then loops back to plan with the answer. A retry
limit stops an endless back-and-forth. \emph{Example:} ``I want you to cook fish''
gets a question, not a made-up plan; a request with six goals gets a question about
\emph{order of work} rather than a guess at the priority.

\section{Rule-Based Reflection}
\label{sec:reflect}
The first plan is just a guess; the system checks and fixes it at three points, each
with simple rules and no model call. These work alongside the model critic
(Section~\ref{sec:plan}-E): the critic checks the plan \emph{with a model}, while these
check the plan, the data, and the result with cheap, easy-to-test rules.

\textbf{A. Plan check, before running.} Before the first agent runs,
\texttt{validate\_and\_repair\_plan} checks that each agent's inputs will be ready,
treating agents as things that need and produce items (repository, code, bugs,
patch, and so on). Two fixes are always safe and do not change the user's intent: a
patch step with no confirmed bugs gets the bug-finder added in front of it; a
re-test step with no patch gets the patch step added (which in turn pulls in the
bug-finder). A step whose input does not exist is removed instead of faked.
\emph{Example:} \texttt{[kio6, kio7]} on a repository becomes
\texttt{[kio5, kio6, kio7]}; \texttt{[kio3]} (``analyze the repository'') with no
repository on disk becomes \texttt{[]}. Agents with no rule are left alone.

\textbf{B. Data check, between steps.} After each agent finishes,
\texttt{reflect\_on\_result} asks whether its result makes any later step pointless,
and removes it. \emph{Example:} if the bug-finder confirms \textbf{zero} bugs, the
queued patch and re-test steps are dropped, so the workflow does not invent work on
empty input.

\textbf{C. Result check, before finishing.} Before a run is marked done,
\texttt{assess\_step} looks at the final signals. \emph{Example:} if the re-test
step reports any failing test, the run is marked \texttt{TESTS\_FAILING} and the
finish step says ``completed with issues'' instead of success---so the report never
claims a passing build when tests failed.

Each check writes a log event, and the plan-fix check uses the same length limit
that keeps the report (Section~\ref{sec:resil}-C), so fixing the plan never quietly drops
the final output.

\section{The Fallback Layer}
\label{sec:resil}
The plan stage is wrapped so that no single failure stops a workflow.

\textbf{A. Model fallback chain.} Plan-stage model calls (planning, clarification,
critic) all go through one place that tries, in order: a strong hosted model when
set up; a local model run inside the cluster; and a fixed default list of steps.
Each step down the chain happens on an error, a rate-limit reply, or an empty
answer, and is logged. Because planning is low-volume, not time-critical, and
returns short JSON, a strong hosted model is cheap here and far more reliable at
valid JSON than the local model---but the system never \emph{depends} on it.
\emph{Example:} when a burst of requests passes the hosted model's rate limit,
planning falls back to the local model and keeps going.

\textbf{B. Facts that fail safely.} Every rule-based fact step (intent, scan,
bidding) and the clarification check returns an empty result on any internal error,
so one flaky check can never block a valid request. Staying up is the default, not a
special case.

\textbf{C. Never silently lose a step.} Length limits that protect against model
mistakes must not delete needed steps. The plan cleaner de-duplicates, applies the
length limit, and---if the limit would drop the \emph{report} agent the planner
chose---puts it back as the last step. The report is the final output; trimming may
never remove it.

\textbf{D. Pausing for a human.} Two automatic signals send the run to a person. An
\emph{unclear request} (the clarification check asked a question) goes to a
checkpoint that asks with options and loops back. A \emph{risky plan} (the planner
fell back, or the critic was \texttt{uncertain}) goes to a review checkpoint that
pauses before the first agent runs, so a bad plan cannot act until a human approves
it.

\textbf{E. Non-blocking warnings.} An agent may notice something it was not asked to
fix (for example, a security concern seen while writing tests). It raises a small,
low-priority \emph{warning} instead of stopping the run; the orchestrator collects
these and reports them with the final result---an extra pair of eyes without the
cost and noise of running every agent on every request.

\section{Evaluation}
\label{sec:eval}
\textbf{A. Method.} Comparing a plan to one ``correct'' plan is a poor test for
large or unclear requests, where many plans are valid. So we test the plan stage by
\emph{behavior}. A test case says, for a request: should it ask for clarification;
which agents \emph{must} appear; which must \emph{not}; and whether the report must
be kept. The harness runs the ask-then-plan chain over a labeled set and scores each
part; the planner is plugged in, so the same test runs against the live hosted model
or a fixed stub. We keep a second, exact-match test for small, clear routing cases.

\textbf{B. What the planner does on a big request.} On a 600-word request with six
goals (analyze, secure, patch, test, build new, optimize, report; with an ``ask if
unclear'' note and an off-topic distractor), the stage: labeled the intent as
security and high risk; asked a clear, multi-option question about the order of
work, and ignored the distractor; when it planned, produced a complete list that
covered every goal and ended with the report; and the critic added the extra
test-automation agent. The hosted model returned valid JSON---including the
left-out list---on the first try, while the small local model used to need the
repair path.

\textbf{C. Case: silently dropping the report.} An early length limit of eight
agents silently dropped the report on the six-goal plan even though the model's own
reason mentioned it. The limit was in \emph{two} places---the planner and the
rule-based plan-fix step (Section~\ref{sec:reflect}-A)---so fixing one was undone by the
other; the behavioral test's report-kept check caught the end-to-end loss. The fix
raised both limits to the number of available agents and added the
put-the-report-back rule (Section~\ref{sec:resil}-C). Re-running confirmed that both the
report and the low-power optimization step survived---exactly the bug an exact-match
test misses and a behavioral test catches.

\textbf{D. Case: a model rate limit.} The live behavioral test found a real limit:
our hosted account allows five requests per minute, while the plan stage makes up to
three hosted calls per workflow (clarify, plan, critic). A burst uses up the budget
and forces the chain down to the local model and then the default list, which the
test correctly scored as degraded. The chain made the system slow down gracefully
instead of failing; the finding points to spacing out requests and planning
capacity.

\textbf{E. Cost and speed.} Because planning returns short output and is low-volume,
the hosted-model cost is a fraction of a cent per plan (less with prompt caching),
and the plan stage takes tens of seconds (clarify, plan, critic). When the request
is unclear, the early question avoids paying for all three calls.

\section{Discussion and Limitations}
\textbf{Reliability comes from the design, not just the model.} Two of the four
problems (silently losing steps, one point of failure) cannot be fixed by a better
model; they are fixed by the design---limits with a put-back rule, fact steps that
fail safely, and the fallback chain. The other two (broken output, wrong choice) are
\emph{reduced} by a stronger model but still need the JSON backup reader and the
clarification gate.

\textbf{Limitations.} Our evaluation is a method plus a few worked cases on a small
labeled set, not a large benchmark, and the behavioral checks encode our judgment of
correct behavior. The critic and the clarification check are themselves model calls
and can be wrong; we reduce the risk with confidence gating and human checkpoints but
do not remove it. The agent fit score is simple word overlap---easy to extend but not
learned. We report one hosted and one local model; others may change the balance
between the JSON backup and the model's own reliability.

\textbf{Future work.} Forcing valid output with a strict schema (to drop the JSON
backup reader); letting an agent decide for itself when it fits a task; spacing out
requests to stay under hosted rate limits; and a larger shared test set built from
real traffic.

\section{Conclusion}
We presented a planner that does not trust one model call to get the plan right.
Instead it builds the plan in small, checkable steps: cheap rule-based facts guide
the choice, a hosted model proposes a plan and says what it left out, a second model
call checks it, and an ``ask, don't guess'' gate handles unclear requests---all
wrapped in a fallback chain with safe-failing checks and human pauses. A behavioral
testing method, rather than exact-match scoring, found real bugs, including a
silently dropped report and a model rate limit. The lesson is that dependable
LLM-driven planning comes less from a bigger model and more from gathering facts,
checking the result, and failing safely.

% Ordered by first citation in the text (IEEE convention), so in-text numbers
% appear as [1], [2], [3], ... in reading order.
\begin{thebibliography}{99}
\bibitem{cot} J. Wei \emph{et al.}, ``Chain-of-Thought Prompting Elicits Reasoning in Large Language Models,'' in \emph{Proc. NeurIPS}, 2022.
\bibitem{react} S. Yao \emph{et al.}, ``ReAct: Synergizing Reasoning and Acting in Language Models,'' in \emph{Proc. ICLR}, 2023.
\bibitem{toolformer} T. Schick \emph{et al.}, ``Toolformer: Language Models Can Teach Themselves to Use Tools,'' in \emph{Proc. NeurIPS}, 2023.
\bibitem{voyager} G. Wang \emph{et al.}, ``Voyager: An Open-Ended Embodied Agent with Large Language Models,'' arXiv:2305.16291, 2023.
\bibitem{autogen} Q. Wu \emph{et al.}, ``AutoGen: Enabling Next-Gen LLM Applications via Multi-Agent Conversation,'' arXiv:2308.08155, 2023.
\bibitem{langgraph} LangChain, ``LangGraph: Stateful, Multi-Actor Applications with LLMs,'' 2024. [Online]. Available: \url{https://langchain-ai.github.io/langgraph/}
\bibitem{reflexion} N. Shinn \emph{et al.}, ``Reflexion: Language Agents with Verbal Reinforcement Learning,'' in \emph{Proc. NeurIPS}, 2023.
\bibitem{judge} L. Zheng \emph{et al.}, ``Judging LLM-as-a-Judge with MT-Bench and Chatbot Arena,'' in \emph{Proc. NeurIPS}, 2023.
\end{thebibliography}

\end{document}
